
\subsection{排列组合的经典方法}

\subsubsection{插空法}

\begin{example}
    5 名男生， 2 名女生，站成一排照相．
    \begin{enumerate}[(1)]
        \item 共有多少种排法？ \underline{\hspace{3cm}}
        \item 两名女生不排在队伍两头的排法有多少种？ \underline{\hspace{3cm}}
        \item 两名女生不相邻的排法有多少种？ \underline{\hspace{3cm}}
    \end{enumerate}
\end{example}

\subsubsection{捆绑法（相邻问题）}
将相邻元素看成一个整体，当作一个元素，参与排列，然后再对相邻元素进行排列。

\begin{example}
    某班优秀学习小组有甲、乙、丙、丁、戊共5人，他们排成一排照相，则甲、乙二人相邻的排法种数为 ( )
    \begin{tasks}(4)
        \task 24
        \task 36
        \task 48
        \task 60
    \end{tasks}
\end{example}

\begin{solution}
    先安排甲、乙相邻，有 $\mathrm{A}_{2}^{2}$ 种排法，再把甲、乙看作一个元素，与其余三个人全排列，故有排法种数为 $\mathrm{A}_{2}^{2} \times \mathrm{A}_{4}^{4} = 2 \times 24 = 48$。故选：C
\end{solution}

\begin{example}
    7人站成一排，其中甲乙相邻且丙丁相邻，共有多少种不同的排法。
\end{example}

\begin{solution}
    可将甲和乙、丙和丁分别捆绑在一起看成两个元素，这两个元素再与其余三人 ($7-2-2=3$人) 一起进行全排列。内部甲乙可以交换，丙丁可以交换。

    总排法数：将 (甲乙) 看作元素X，(丙丁) 看作元素Y。现在有 X, Y, 和另外3个人，共5个元素排列，有 $\mathrm{A}_{5}^{5}$ 种。

    甲乙内部排列有 $\mathrm{A}_{2}^{2}$ 种。
    丙丁内部排列有 $\mathrm{A}_{2}^{2}$ 种。

    故共有 $\mathrm{A}_{5}^{5} \mathrm{A}_{2}^{2} \mathrm{A}_{2}^{2} = 120 \times 2 \times 2 = 480$ 种不同的排列方法。
\end{solution}

\subsubsection{插空法（不相邻问题）}
元素不相邻问题，可先把无位置要求的几个元素全排列，再把规定相离的几个元素插入上述几个元素间的空位（包含两端）。

\begin{example}
    省实验中学为预防秋季流感爆发，计划安排学生在校内进行常规体检，共有3个检查项目，需要安排在3间空教室进行检查，学校现有高一排6间的空教室供选择使用，但是为了避免学生拥挤，要求作为检查项目的教室不能相邻，则共有 (\hspace{0.7cm}
    ) 种安排方式。
    \begin{tasks}(4)
        \task 12
        \task 24
        \task 36
        \task 48
    \end{tasks}
\end{example}

\begin{solution}
    6间空教室，选3间用于检查且不相邻。可以看作3间被选中的教室 (C) 和3间未被选中的空教室 (E)。
    为了使选中的教室不相邻，先排列3间未被选中的空教室 E E E。这会产生4个空位（包括两端）：\_ E \_ E \_ E \_。
    将3间被选中的教室安排在这4个空位中，每个空位至多一间教室，有 $\mathrm{A}_{4}^{3}$ 种安排方式。
    $\mathrm{A}_{4}^{3} = 4 \times 3 \times 2 = 24$。
    故选：B。
\end{solution}

\begin{example}
    3 位老师和 4 名学生站成一排，要求任意两位老师都不相邻，则不同的排法种数为 (\hspace{0.7cm})
    \begin{tasks}(4)
        \task $\mathrm{A}_{4}^{4}$
        \task $\mathrm{A}_{5}^{3} \mathrm{A}_{4}^{4}$
        \task $\mathrm{A}_{7}^{7}$
        \task $\mathrm{A}_{4}^{4} \mathrm{A}_{3}^{3}$
    \end{tasks}
\end{example}

\begin{solution}
    根据题意，分2步进行：

    \begin{enumerate}
        \item 将4名学生站成一排，有 $\mathrm{A}_{4}^{4}$ 种排法；
        \item 4名学生排好后，形成5个空位（包括两端）：\_ S \_ S \_ S \_ S \_。
              将3位老师插入这5个空位中，每位老师占一个空位，有 $\mathrm{A}_{5}^{3}$ 种排法。
    \end{enumerate}

    则总的排法种数为 $\mathrm{A}_{4}^{4} \times \mathrm{A}_{5}^{3} = 24 \times 60 = 1440$ 种。
\end{solution}

\subsubsection{平均分组问题}

\begin{example}
    6本不同书平均分成3堆，每堆2本共有多少种分法？
\end{example}

\begin{solution}
    解：分三步取书，每次取两本。
    第一堆从6本中选2本，有 $\mathrm{C}_{6}^{2}$ 种方法。
    第二堆从剩下的4本中选2本，有 $\mathrm{C}_{4}^{2}$ 种方法。
    第三堆从剩下的2本中选2本，有 $\mathrm{C}_{2}^{2}$ 种方法。

    因为三堆书的数量相同（都是2本），所以堆与堆之间没有顺序，属于无序分组。若按上述步骤计算 $\mathrm{C}_{6}^{2} \mathrm{C}_{4}^{2} \mathrm{C}_{2}^{2}$，则对堆进行了排序。

    例如，若6本书为A,B,C,D,E,F。一种分法是 (AB), (CD), (EF)。在上述计算中，(AB, CD, EF), (AB, EF, CD), (CD, AB, EF), (CD, EF, AB), (EF, AB, CD), (EF, CD, AB) 这 $\mathrm{A}_{3}^{3}=6$ 种情况都被算作不同的情况，但实际上它们是同一种分组结果。

    因此，需要除以堆数的阶乘 $\mathrm{A}_{3}^{3}$ (或者 $3!$)。
    所以共有 $$ \frac{\mathrm{C}_{6}^{2} \mathrm{C}_{4}^{2} \mathrm{C}_{2}^{2}}{\mathrm{A}_{3}^{3}} = \frac{15 \times 6 \times 1}{6} = 15$$ 种分法。
\end{solution}

\subsubsection{分组分配问题}

先分组，再分配，分组时如果出现了均分的情况，需要除以相同组数的全排列。下面是一道母题，掌握后可以杀所有这样的题。

\vspace{0.4cm}

\begin{example}
    将 4 名优秀教师分配到 3 个不同的学校进行教学交流，每名优秀教师只分配到 1 个学校，每个学校至少分配 1 名优秀教师，则不同的分配方案共有
    \begin{tasks}(4)
        \task $72$ 种
        \task $48$ 种
        \task $36$ 种
        \task $24$ 种
    \end{tasks}
\end{example}

\begin{solution}
    首先，分配四名教师到三个学校，只有以下这种情况：

    $$ 4 = 1 + 1 + 2.$$

    我们视为从 4 个人中选出 1 个，则为 $C_4^1$，那么还剩下 3 个人，再选出 1 个，则为 $C_3^1$，剩下的就是 $C_2^2$，这个过程是分步骤的，根据分步乘法原理，应该最后结果是 $C_4^1C_3^1C_2^2$ 种选人的方式。

    重点！这里分出的 $1 + 1$ 实际上是有重合的情况的，所以我们要除以【相同人数组数的全排列】，即 $\frac{C_4^1C_3^1C_2^2}{A_2^2}$，$\frac{C_4^1C_3^1C_2^2}{A_2^2}$ 就相当于我们分组的方案数了，再用这些去进行全排列即可 ($\times A_3^3$)
\end{solution}

\begin{problem}
如果是 $5$ 名呢？思考一下。
\end{problem}

\begin{example}
    (2020·全国) 疫情防控期间，上海某医院安排5名专家到3个不同的区级医院支援，每名专家只去一个区级医院，每个区级医院至少安排一名专家，则不同的安排方法共有 (\hspace{0.6cm})
    \begin{tasks}(4)
        \task 60种
        \task 90种
        \task 150种
        \task 240种
    \end{tasks}
\end{example}

\subsubsection{分组分配练习题}

\textbf{下面的练习题难度逐渐递增}。
\vspace{0.4cm}

\begin{enumerate}
    \item 为提升学生综合素养，某高中开设了＂围棋＂、＂机器人＂、＂乒乓球＂、＂建筑设计＂四门选修课程，要求每位学生每学年至多选两门课程，高一至高三三学年必须修完四门课程，每学年上学期选一次，每门课程限选修一学年，则每位学生不同的选修方式有
          \begin{tasks}(4)
              \task $9$ 种
              \task $36$ 种
              \task $54$ 种
              \task $72$ 种
          \end{tasks}


    \item 现安排甲、乙、丙、丁、戊 5 位志愿者到三个社区做志愿服务工作，每个社区至少安排1人，每位志愿者只到一个社区，其中甲、乙安排在同一个社区的概率为
          \begin{tasks}(4)
              \task $\frac{3}{25}$
              \task $\frac{1}{5}$
              \task $\frac{6}{25}$
              \task $\frac{3}{10}$
          \end{tasks}

    \item 某校组织校运会活动，由甲、乙、丙三名志愿者负责 $A, B, C, D$ 四个任务，每人至少负责一个任务，每个任务都有且仅有一人负责，且甲不负责 A 任务，则不同的任务分配方法种数为
          \begin{tasks}(4)
              \task $12$
              \task $18$
              \task $24$
              \task $30$
          \end{tasks}

    \item 光正实验学校高二年级拟举行＂诗词＂、＂历史＂、＂地理＂三场不同主题的知识竞答活动，要求各班各派 3 名学生分别参加这三个主题的竞答。某班准备从甲、乙、丙、丁 4 位同学中选派 3 位，已知甲不参加＂诗词＂主题的竞答活动，则该班不同的选派方法有
          \begin{tasks}(4)
              \task $9$ 种
              \task $12$ 种
              \task $15$ 种
              \task $18$ 种
          \end{tasks}

    \item 从 5 名同学中选择 4 人参加三天志愿服务活动，有一天安排两人，另两天各安排一人，共有 \underline{\hspace{2cm}}种安排方法（用数字作答）

    \item 总院计划安排小吴、小张等 8 名护士到 $A, B, C, D$ 四所医院实习，每两名护士到一所医院，考虑到个人意愿问题，小吴不安排 A 医院，小张不安排 $B$ 医院，则安排方法共有
          \begin{tasks}(4)
              \task $1200$ 种
              \task $1440$ 种
              \task $1600$ 种
              \task $1800$ 种
          \end{tasks}

\end{enumerate}

\begin{example}
    有7个不同的社区需要3名志愿者（甲乙丙）都去参加志愿服务，要求每个社区至多安排一人，每名志愿者恰好去一个社区，则不同的安排方法有多少种？
\end{example}

\subsubsection{5、特殊元素或位置优先考虑}
\begin{example}
    由0,1,2,3,4,5可以组成多少个没有重复数字的五位奇数。
\end{example}
\begin{solution}
    解：一个五位奇数，需要满足：
    \begin{enumerate}
        \item 末位是奇数：从 \{1, 3, 5\} 中选一个，有 $\mathrm{C}_{3}^{1}$ 种选法。
        \item 首位不能是0，且不能是末位已选的数字。
        \item 中间三位从剩余数字中选取并排列。
    \end{enumerate}
    采用优先安排特殊位置的方法：
    \begin{enumerate}
        \item 排末位 (奇数)：从 \{1, 3, 5\} 中选1个数字，有 $\mathrm{C}_{3}^{1} = 3$ 种选择。
        \item 排首位 (非0且非末位数字)：
              可供选择的数字有6个 \{0,1,2,3,4,5\}。去掉末位已选的1个数字，剩下5个。
              如果0未被末位选中：
              \begin{itemize}
                  \item 若0是可选的，则首位不能为0，也不能为末位数字。所以有 $6-1(\text{末位})-1(0) = 4$ 种选择。
              \end{itemize}
        \item 排中间三位：从剩下的 $6-2=4$ 个数字中选3个进行全排列，有 $\mathrm{A}_{4}^{3}$ 种方法。
    \end{enumerate}
    组合计算：
    先选末位奇数：$\mathrm{C}_{3}^{1}$ 种。
    假设末位选了一个奇数 (例如1)。剩下 \{0,2,3,4,5\}。
    选首位：首位不能是0，所以从 \{2,3,4,5\} 中选一个，有 $\mathrm{C}_{4}^{1}$ 种。
    假设首位选了一个 (例如2)。剩下 \{0,3,4,5\} (4个数字)。
    选中间三位：从这4个数字中选3个排列，有 $\mathrm{A}_{4}^{3}$ 种。
    所以总数为：$\mathrm{C}_{3}^{1} \times \mathrm{C}_{4}^{1} \times \mathrm{A}_{4}^{3} = 3 \times 4 \times (4 \times 3 \times 2) = 3 \times 4 \times 24 = 12 \times 24 = 288$ 个。
\end{solution}

\subsubsection{定序问题\*也称受限区域排列}
所谓定序问题是指几个元素顺序固定的排列问题，解法是先将这几个元素与其它元素一同进行全排列，然后除以定序元素的全排列数。或者，先为定序元素选定位置，其余元素进行全排列。

\begin{example}
    例7、将数字1，2，3，4填入标号为1，2，3，4的四个方格里，每格填一个数，则每个方格的标号与所填数字均不相同的填法有（ ）
    \begin{tasks}(4)
        \task 6
        \task 9
        \task 11
        \task 23
    \end{tasks}
\end{example}
\begin{solution}
    解：此题为“错位排列”问题。记 $D_n$ 为n个元素的错位排列数。
    $D_1 = 0$
    $D_2 = 1$ (21)
    $D_3 = 2$ (231, 312)
    $D_4 = 9$
    递推公式为 $D_n = (n-1)(D_{n-1} + D_{n-2})$ for $n \ge 2$。
    $D_4 = (4-1)(D_3 + D_2) = 3 \times (2+1) = 3 \times 3 = 9$。
    另一种公式是 $D_n = n! \left(1 - \frac{1}{1!} + \frac{1}{2!} - \frac{1}{3!} + \dots + (-1)^n \frac{1}{n!}\right)$。
    $D_4 = 4! \left(\frac{1}{2!} - \frac{1}{3!} + \frac{1}{4!}\right) = 24 \left(\frac{1}{2} - \frac{1}{6} + \frac{1}{24}\right) = 12 - 4 + 1 = 9$。
    对于本题，4个元素错位排列，共有9种。
    故选B。
\end{solution}

\noindent\textbf{元素排列满足一定的要求，分类讨论法}
\newline
解题时，常常需要根据题意，把所有可能发生的情况，分成若干类，然后逐类分别进行研究，最后把各类结果汇总，得出问题的解答。分类的原则是不重不漏。

\begin{example}
    由数字0，1，2，3，4，5组成没有重复数字的六位数，其中个位数字小于十位数字的共有（ ）个。
    \begin{tasks}(4)
        \task 210
        \task 300
        \task 464
        \task 600
    \end{tasks}
\end{example}
\begin{solution}
    解：考虑总的没有重复数字的六位数，再考虑个位与十位的关系。
    方法：先从6个数字中选6个进行排列，首位不能为0。总的六位数有 $\mathrm{A}_{6}^{6} - \mathrm{A}_{5}^{5}$ (如果0在首位) $= 5 \times \mathrm{A}_{5}^{5} = 5 \times 120 = 600$ 个。
    对于任意选定的两个位置（个位和十位）和任意选定的两个不同数字，这两个数字在这两个位置只有一种排列方式使得个位小于十位。
    考虑后两位（个位和十位）。从6个数字中选2个数字，有 $\mathrm{C}_{6}^{2}$ 种选法。对于选定的2个数字，只有1种方式满足个位<十位。
    剩下的4个数字排前四位，首位不能为0。

    Alternative approach (as per image solution):
    按最高位是否为0分类（这里是组成六位数，所以最高位不能为0）。
    \begin{enumerate}
        \item 选首位：从 \{1,2,3,4,5\} 中选1个，有 $\mathrm{C}_{5}^{1}$ 种 (或 $\mathrm{A}_{5}^{1}$)。
        \item 选个位和十位：从剩下的5个数字中选2个数字，有 $\mathrm{C}_{5}^{2}$ 种选法。由于要求个位数字小于十位数字，这两个选出的数字在个位和十位的排法是固定的 (1种)。
        \item 排其余三位 (十万位已定，个位十位数字已定，剩下3个位置)：从剩下的 $6-1(\text{首位})-2(\text{个/十位})=3$ 个数字中选3个全排列，有 $\mathrm{A}_{3}^{3}$ 种。
    \end{enumerate}
    总共 $C_{5}^{1} \times C_{5}^{2} \times 1 \times A_{3}^{3} = 5 \times 10 \times 1 \times 6 = 300$ 个。
    故选B。
\end{solution}

\subsubsection{元素相同问题（隔板法）}
将 $n$ 个相同元素分成 $m$ 份（$m$,$n$ 为正整数，每份至少一个元素），可以利用 $m-1$ 块隔板，插入 $n$ 个元素所形成的 $n-1$ 个空隙（不含两端）中的任意 $m-1$ 个空位，共有 $\mathrm{C}_{n-1}^{m-1}$ 种分法。

\begin{example}
    有10个运动员名额，要分给7个班，每班至少一个，有多少种分配方案？
\end{example}
\begin{solution}
    解：此问题看成将10个相同的名额（元素）分配给7个不同的班级，每个班级至少一个名额。
    这等价于将10个相同的小球放入7个不同的盒子，每个盒子至少一个小球。
    使用隔板法：10个名额看作10个小球 O O O O O O O O O O。
    它们之间有 $10-1=9$ 个空隙。
    要分成7份，需要插入 $7-1=6$ 块隔板。
    从9个空隙中选择6个位置放隔板，有 $\mathrm{C}_{9}^{6}$ 种方法。
    $\mathrm{C}_{9}^{6} = \mathrm{C}_{9}^{9-6} = \mathrm{C}_{9}^{3} = \frac{9 \times 8 \times 7}{3 \times 2 \times 1} = 3 \times 4 \times 7 = 84$ 种。
\end{solution}

\subsubsection{交叉问题（容斥原理）}
某些排列组合问题，在计算时会出现重复计数的情况，直接计算比较困难，这时可以利用容斥原理来解决。两个集合的容斥原理：

$$card(A \cup B) = card(A) + card(B) - card(A \cap B)$$

对于三个集合:

\begin{eqnarray*}
    card(A \cup B \cup C) &=& card(A) + card(B) + card(C) \\
    &&- card(A \cap B) - card(A \cap C) - card(B \cap C) \\
    &&+ card(A \cap B \cap C)
\end{eqnarray*}

\begin{example}
    从6名运动员中选出4人参加 $4 \times 100$米接力赛，如果甲不跑第一棒，乙不跑最后一棒，共有多少种不同的参赛方法？
\end{example}
\begin{solution}
    解：设全集 $U$ 为从6名运动员中选出4人并排列参加接力赛的总方法数。
    $|U| = \mathrm{A}_{6}^{4} = 6 \times 5 \times 4 \times 3 = 360$ 种。

    设 $A$ 为甲跑第一棒的事件（或排法集合）。
    如果甲跑第一棒，则需要在剩下5人中选3人跑其余三棒并排列：$|A| = 1 \times \mathrm{A}_{5}^{3} = 5 \times 4 \times 3 = 60$ 种。

    设 $B$ 为乙跑最后一棒的事件。
    如果乙跑最后一棒，则需要在剩下5人中选3人跑前三棒并排列：$|B| = \mathrm{A}_{5}^{3} \times 1 = 60$ 种。

    $A \cap B$ 为甲跑第一棒且乙跑最后一棒的事件。
    甲跑第一棒，乙跑最后一棒。则需要在剩下 $6-2=4$ 人中选2人跑中间两棒并排列：$|A \cap B| = 1 \times \mathrm{A}_{4}^{2} \times 1 = 4 \times 3 = 12$ 种。

    所求为甲不跑第一棒且乙不跑最后一棒，即求 $|U| - |A \cup B|$。
    $|A \cup B| = |A| + |B| - |A \cap B| = 60 + 60 - 12 = 108$ 种。
    所以，符合条件的排法数 $= |U| - |A \cup B| = 360 - 108 = 252$ 种。
\end{solution}

\subsubsection{正难则反思想}

\begin{example}
    有五名学生站成一排，从中任选三人，要求这三人中任何两人都不相邻有多少种选法？
\end{example}

\begin{solution}
    设 $5$ 名学生的位置为 1, 2, 3, 4, 5。要选出3个位置，使得这3个位置不相邻。
    符合条件的选择只有一种：选位置 1, 3, 5。
    所以有 $\mathrm{C}_{1}^{1} = 1$ 种选法。(如果学生是无差别的，只考虑位置选法)
    如果学生是不同的，那么选出这三名学生后，他们站在1,3,5号位。
    题目问的是“选法”，通常指选择学生组合或位置组合。
    使用公式法：从n个元素中选k个不相邻的元素，有 $\mathrm{C}_{n-k+1}^{k}$ 种选法。
    这里 $n=5$ (五个位置)，$k=3$ (选三个人/位置)。
    种数 $= \mathrm{C}_{5-3+1}^{3} = \mathrm{C}_{3}^{3} = 1$ 种。
\end{solution}

\subsubsection{选排问题}
\begin{example}
    9名乒乓球运动员中，其中男运动员5名，女运动员4名，现要进行混合双打比赛（即每队由一名男运动员和一名女运动员组成，共组成两队进行比赛），有多少种不同的配对方法？
\end{example}
\begin{solution}
    解：要组成两对混合双打队伍。
    步骤1：从5名男运动员中选出2名男运动员。有 $\mathrm{C}_{5}^{2}$ 种方法。
    步骤2：从4名女运动员中选出2名女运动员。有 $\mathrm{C}_{4}^{2}$ 种方法。
    假设选出的男运动员为 M1, M2；女运动员为 F1, F2。
    步骤3：将这2男2女组成两对混合双打。
    M1可以和F1配对，则M2和F2配对：(M1,F1) 和 (M2,F2)。
    或者M1可以和F2配对，则M2和F1配对：(M1,F2) 和 (M2,F1)。
    有2种配对方式。
    所以总的配对方法数 = $\mathrm{C}_{5}^{2} \times \mathrm{C}_{4}^{2} \times 2 = 10 \times 6 \times 2 = 120$ 种。
\end{solution}

\subsubsection{有序排列问题}

\begin{example}
    某学校举行校园歌手大赛活动邀请了 6 位专家评委，在活动结束时邀请这 6 位专家站成一排合影留念，则下列说法正确的是
    \begin{tasks}(1) % Options are long, better as single column or (2)
        \task 若将专家甲、乙、丙三人从左到右按照身高递增的顺序排列，则共有 120 种排法
        \task 若专家甲、乙两人不相邻，则共有 480 种排法
        \task 若专家甲、乙、丙三人相邻，且甲在中间，则共有 72 种排法
        \task 若专家丙不在排头，丁不在排尾，则共有 480 种排法
    \end{tasks}
\end{example}

\begin{enumerate}
    \item 某班一天上午有 4 节课，下午有 2 节课，现要安排该班一天中语文、数学、政治、英语、体育、心理 6 堂课的课程表，要求数学课排在上午，体育课排在下午，不同排法总数是
          \begin{tasks}(4)
              \task $192$
              \task $144$
              \task $124$
              \task $216$
          \end{tasks}
\end{enumerate}
